\documentclass[10pt,letterpaper]{article}
\usepackage{cornell-notes}

\title{Chapter 1: Introduction}
\author{}
\date{}

\setDocTitle{Chapter 1: Introduction}
\setDocSubtitle{Computer Networks}
\setDocVersion{v1.0}
\setDocStatus{Study Companion}
\setDocOwner{Computer Networks Cornell Notes}
\setDocAuthor{}
\setDocDate{\today}
\setDocSummary{Cornell-format study and review notes for Chapter 1: Introduction.}

\setCornellCollection{Computer Networks Cornell Notes}
\setCornellUnitType{Chapter}
\setCornellUnitNumber{1}
\setCornellUnitTitle{Introduction}
\setCornellCompanionLabel{Study and review companion}

\hypersetup{pdftitle={Chapter 1: Introduction},pdfsubject={Cornell notes},pdfauthor={}}

\begin{document}
\maketitle
\notesection{Chapter overview}
\begin{overviewbox}
\textbf{Learning objectives}\begin{itemize}
  \item Relate business, home, mobile, and social uses to network requirements.
  \item Classify networks by scale and distinguish a network from an internetwork.
  \item Explain protocol layering, services, interfaces, and connection models.
  \item Compare the OSI and TCP/IP reference models without confusing a model with an implementation.
  \item Place the Internet, cellular systems, wireless LANs, RFID, and sensor networks in a common conceptual framework.
\end{itemize}
\textbf{Prerequisite concepts}\quad Basic computer organization, binary data, and client/server terminology.\par
\textbf{Key themes}\quad Resource sharing; communication; scale; layered abstraction; protocol/service separation; interoperability; standardization.\par
\textbf{Named mechanisms and standards}\quad OSI model, TCP/IP model, Internet architecture, third-generation cellular architecture, IEEE 802.11, RFID, sensor networks, ITU-T, ISO, IEEE, and IETF processes.
\end{overviewbox}

\notesection{Cornell cue-and-note record}

\begin{longtable}{@{}L{0.245\textwidth}|L{0.695\textwidth}@{}}
\rowcolor{CNavy}\color{white}\bfseries Cue / recall prompt &
\color{white}\bfseries Notes \\
\endfirsthead
\rowcolor{CNavy}\color{white}\bfseries Cue / recall prompt &
\color{white}\bfseries Notes (continued) \\
\endhead
\hline
\endfoot
\cellcolor{CCue}\textbf{1.1 Uses of computer networks} & \begin{minipage}[t]{\linewidth}\raggedright Network value comes from applications and shared access; each use imposes different capacity, delay, reliability, mobility, and security needs.\par\begin{itemize}\item \textbf{1.1.1 Business applications} Networks share information and equipment, support client-server and peer communication, and enable electronic commerce. A server provides a service; a client requests it. The requirement is dependable access to distributed resources, not merely physical connectivity.\item \textbf{1.1.2 Home applications} Residential networks support information access, person-to-person communication, interactive entertainment, and electronic commerce. Peer-to-peer designs distribute roles among end systems instead of requiring one permanent server.\item \textbf{1.1.3 Mobile users} Portable devices combine mobility with wireless access. Wireless and mobile are distinct: a stationary host may use radio, while a mobile host may retain service through changing points of attachment.\item \textbf{1.1.4 Social issues} Connectivity raises questions about speech, filtering, privacy, monitoring, ownership, and acceptable behavior. Technical mechanisms can enforce policy, but the policy itself is a social and legal choice.\end{itemize}\end{minipage} \\ \hline
\cellcolor{CCue}\textbf{1.2 Network hardware} & \begin{minipage}[t]{\linewidth}\raggedright Scale, ownership, and transmission organization provide useful classifications. Broadcast media share one channel; point-to-point systems select paths through intermediate nodes.\par\begin{itemize}\item \textbf{1.2.1 Personal area networks} A personal area network connects devices around one person over a short range; cable and short-range radio are typical media.\item \textbf{1.2.2 Local area networks} A LAN covers a home, office, or campus-scale site. Wired Ethernet commonly uses switched point-to-point links; wireless LANs use a shared radio channel.\item \textbf{1.2.3 Metropolitan area networks} A MAN spans a city-scale area and often joins access networks with provider infrastructure.\item \textbf{1.2.4 Wide area networks} A WAN spans large geographic regions. Hosts attach through access links to a communication subnet whose routers store and forward packets over transmission lines.\item \textbf{1.2.5 Internetworks} An internetwork joins unlike networks through devices that reconcile their boundaries. The Internet is an internetwork, whereas a single technology deployed broadly is still one network.\end{itemize}\end{minipage} \\ \hline
\cellcolor{CCue}\textbf{1.3 Network software} & \begin{minipage}[t]{\linewidth}\raggedright The central discipline is to separate externally visible service behavior from the internal peer protocol that realizes it.\par\begin{itemize}\item \textbf{1.3.1 Protocol hierarchies} Layers reduce complexity. Peer entities follow a protocol; each layer offers services upward and uses the service below. Headers carry control information, while encapsulation nests higher-layer data inside lower-layer protocol data units.\item \textbf{1.3.2 Design issues for the layers} Recurring issues include addressing, error control, flow control, multiplexing, routing, ordering, and connection establishment. The layer chosen for a mechanism affects both scope and cost.\item \textbf{1.3.3 Connection-oriented versus connectionless service} Connection-oriented service establishes state, transfers data, and releases the association; connectionless service treats each message independently. Either can be reliable or unreliable.\item \textbf{1.3.4 Service primitives} Primitives such as LISTEN, CONNECT, ACCEPT, RECEIVE, SEND, and DISCONNECT describe operations visible at an interface, not the messages exchanged between peer protocols.\item \textbf{1.3.5 Services versus protocols} A service specifies what a layer offers to its users. A protocol specifies how peer entities cooperate to provide it. An interface specifies how a user accesses the service.\end{itemize}\end{minipage} \\ \hline
\cellcolor{CCue}\textbf{1.4 Reference models} & \begin{minipage}[t]{\linewidth}\raggedright A reference model organizes responsibilities; it is not itself a complete protocol suite.\par\begin{itemize}\item \textbf{1.4.1 OSI reference model} Seven layers separate physical transmission, data-link framing, network routing, transport, session, presentation, and application functions. The model emphasizes explicit services, interfaces, and protocols.\item \textbf{1.4.2 TCP/IP reference model} The Internet stack groups link access below an internetwork layer, supports end-to-end transport, and places application protocols above transport. IP is the packet-delivery foundation; TCP and UDP provide different transport services.\item \textbf{1.4.3 Model used for study} A five-layer teaching model uses physical, data link, network, transport, and application layers, while security is treated across layers.\item \textbf{1.4.4 Comparing OSI and TCP/IP} Both use layering and transport-level end-to-end communication. OSI clearly distinguishes service, interface, and protocol; TCP/IP grew around deployed protocols and has less sharply separated lower layers.\item \textbf{1.4.5 Critique of OSI} The protocols suffered from timing, technology, implementation, and policy problems; a sound conceptual model did not guarantee successful protocol adoption.\item \textbf{1.4.6 Critique of TCP/IP} Its model does not cleanly separate concepts, is not a general model for arbitrary stacks, and leaves host-to-network behavior underspecified. Its protocols nevertheless achieved wide deployment.\end{itemize}\end{minipage} \\ \hline
\cellcolor{CCue}\textbf{1.5 Example networks} & \begin{minipage}[t]{\linewidth}\raggedright Examples reveal that a common layered vocabulary can describe networks with very different media, administrative structures, and mobility patterns.\par\begin{itemize}\item \textbf{1.5.1 The Internet} End systems connect through access networks to Internet service providers. Routers forward packets across interconnected provider networks; applications operate at the edge while IP supplies common internetwork delivery.\item \textbf{1.5.2 Third-generation mobile networks} Cellular systems divide coverage into cells, authenticate subscribers, route calls and data through operator cores, and hand users between base stations. Circuit and packet domains coexist in the presented architecture.\item \textbf{1.5.3 Wireless LANs: 802.11} Stations associate with an access point in infrastructure mode or communicate directly in ad hoc mode. Shared radio and mobility require link procedures that differ from switched Ethernet.\item \textbf{1.5.4 RFID and sensor networks} RFID readers interrogate tags for identification; sensor networks coordinate many low-power nodes that observe and report physical conditions. Energy and scale strongly constrain protocol design.\end{itemize}\end{minipage} \\ \hline
\cellcolor{CCue}\textbf{1.6 Network standardization} & \begin{minipage}[t]{\linewidth}\raggedright Standards permit independently produced equipment and software to interoperate and also reflect economic and political coordination.\par\begin{itemize}\item \textbf{1.6.1 Telecommunications organizations} National and international carriers coordinate through bodies such as the ITU-T, whose recommendations promote compatible telecommunication services and interfaces.\item \textbf{1.6.2 International standards organizations} ISO and related organizations develop formal international standards; IEEE 802 has major responsibility for LAN and MAN standards.\item \textbf{1.6.3 Internet standards organizations} The Internet community develops specifications through working groups and open documents. Requests for Comments range from proposals and information to standards-track specifications.\end{itemize}\end{minipage} \\ \hline
\cellcolor{CCue}\textbf{1.7 Metric units} & \begin{minipage}[t]{\linewidth}\raggedright Data rates use powers of 1000 (kbit/s, Mbit/s, Gbit/s), while memory sizes are often described with powers of 1024. Always inspect the context before interpreting a prefix.\end{minipage} \\ \hline
\cellcolor{CCue}\textbf{1.8 Structural roadmap} & \begin{minipage}[t]{\linewidth}\raggedright The layered roadmap proceeds upward from signal transmission, through framing and medium access, routing, end-to-end transport, and applications; security mechanisms cut across this sequence.\end{minipage} \\ \hline
\cellcolor{CCue}\textbf{1.9 Summary} & \begin{minipage}[t]{\linewidth}\raggedright Networks connect users and resources across scales. Hardware supplies paths, software layers manage complexity, protocols coordinate peers, services define layer boundaries, and standards make heterogeneous implementations work together.\end{minipage} \\ \hline
\end{longtable}

\begin{legacybox}The third-generation cellular architecture and several historical protocol critiques are preserved as presented. They explain the evolution of networking but should not be mistaken for a current deployment inventory.\end{legacybox}

\notesection{Terminology and notation}

\begin{longtable}{@{}L{0.245\textwidth}|L{0.695\textwidth}@{}}
\rowcolor{CNavy}\color{white}\bfseries Cue / recall prompt &
\color{white}\bfseries Notes \\
\endfirsthead
\rowcolor{CNavy}\color{white}\bfseries Cue / recall prompt &
\color{white}\bfseries Notes (continued) \\
\endhead
\hline
\endfoot
\cellcolor{CCue}\textbf{Protocol} & \begin{minipage}[t]{\linewidth}\raggedright Rules and message conventions used by peer entities at the same layer.\end{minipage} \\ \hline
\cellcolor{CCue}\textbf{Service} & \begin{minipage}[t]{\linewidth}\raggedright Capabilities a lower layer exposes to the layer above, independent of how peers implement them.\end{minipage} \\ \hline
\cellcolor{CCue}\textbf{Interface} & \begin{minipage}[t]{\linewidth}\raggedright Operations and parameters by which one layer accesses another layer's service.\end{minipage} \\ \hline
\cellcolor{CCue}\textbf{Encapsulation} & \begin{minipage}[t]{\linewidth}\raggedright Placement of a higher-layer protocol data unit inside a lower-layer payload with additional control information.\end{minipage} \\ \hline
\cellcolor{CCue}\textbf{Host / end system} & \begin{minipage}[t]{\linewidth}\raggedright A computer or device that runs applications at the edge of a network.\end{minipage} \\ \hline
\cellcolor{CCue}\textbf{Router} & \begin{minipage}[t]{\linewidth}\raggedright A network-layer device that selects outgoing links and forwards packets among networks.\end{minipage} \\ \hline
\cellcolor{CCue}\textbf{LAN / MAN / WAN} & \begin{minipage}[t]{\linewidth}\raggedright Local-, metropolitan-, and wide-area network classifications based primarily on geographic scope and organization.\end{minipage} \\ \hline
\cellcolor{CCue}\textbf{Internetwork} & \begin{minipage}[t]{\linewidth}\raggedright A collection of distinct networks joined so that hosts can communicate across network boundaries.\end{minipage} \\ \hline
\cellcolor{CCue}\textbf{Connection-oriented} & \begin{minipage}[t]{\linewidth}\raggedright Service organized around setup, ordered data transfer, and release, usually with maintained state.\end{minipage} \\ \hline
\cellcolor{CCue}\textbf{Connectionless} & \begin{minipage}[t]{\linewidth}\raggedright Service in which independently addressed messages can be handled without prior setup.\end{minipage} \\ \hline
\end{longtable}

\notesection{Comparison matrix}

\begin{longtable}{@{}L{0.313\textwidth}L{0.313\textwidth}L{0.313\textwidth}@{}}
\toprule
\textbf{Concept} & \textbf{What it specifies} & \textbf{Do not confuse it with} \\ \midrule
Service & Behavior available to an upper-layer user & The peer message exchange \\
Protocol & Rules used between peer entities & The local application programming interface \\
Interface & How adjacent layers invoke service operations & A reference model or wire format \\
Connection-oriented & An association with setup and state & Automatic reliability; a connection can still be unreliable \\
Connectionless & Independent messages without prior setup & Automatic unreliability; acknowledgments may be added \\
\bottomrule
\end{longtable}
\notesection{Chapter synthesis}

\begin{overviewbox}\textbf{Cause-and-effect chain.} Diverse applications create different communication requirements; heterogeneity creates complexity; layering localizes that complexity; protocols coordinate peers; services stabilize layer boundaries; standards allow independent implementations to interoperate.\par
\textbf{Common confusions.} A wireless network is not necessarily mobile; a WAN is not the same as the Internet; OSI is a model rather than a deployed suite; and connection orientation does not by itself imply reliability.\par
\textbf{Derived insight.} Layering is a controlled tradeoff: abstraction improves modularity, but duplicated functions and hidden cross-layer information can reduce efficiency.\end{overviewbox}

\notesection{Retrieval practice}

\begin{enumerate}
  \item How do business and mobile applications impose different requirements on a network?
  \item Distinguish a LAN, a WAN, and an internetwork using scope and structure.
  \item Trace encapsulation as an application message moves down a five-layer stack.
  \item Why are service, protocol, and interface separate concepts?
  \item Can a connectionless service be reliable? Explain without relying on a protocol name.
  \item Compare the seven-layer OSI model with the Internet-oriented model.
  \item Why did a conceptually strong reference model fail to guarantee successful protocol adoption?
  \item Describe the roles of hosts, access networks, providers, and routers in the Internet.
  \item Why does shared radio change link-layer design compared with switched Ethernet?
  \item Contrast RFID and sensor-network goals and constraints.
  \item What problem does standardization solve, and why are multiple standards organizations involved?
  \item Why must a reader distinguish decimal data-rate prefixes from binary storage conventions?
\end{enumerate}
\begin{CornellSummaryBox}{Rapid-review Cornell summary}Applications motivate connectivity; hardware supplies transmission paths; layered software organizes communication; peer protocols realize services; interfaces expose those services; reference models allocate responsibilities; standards enable interoperability. Keep service, protocol, interface, network, and internetwork conceptually distinct.\end{CornellSummaryBox}
\end{document}
